% Level ceiling (grade 5): speeds as "kilometres in one hour" in words,
% whole-number speeds, distance = speed x time with whole numbers only.
% No v = d/t with letters (literal expressions arrive in math grade 7),
% no \unit{km/h} compound notation (arrives grade 7 with quotients).
\chapter{La vitesse~: distance et durée}\label{ch:g5:speed-distance-time}

«~Je suis plus rapide~!~» --- «~Non, c'est moi~!~» Toutes les cours
de récréation tranchent cela honnêtement~: on s'aligne, à vos
marques, prêts, partez. Mais comment comparer un cycliste et un
coureur qui a couru hier, sur une autre route~? La \omterm{def:g5:speed-distance-time:speed}{vitesse} est la
façon dont la physique tranche les courses qui ne pourront jamais
être disputées côte à côte.

\section{Deux façons de gagner une course}

\begin{proposition}[Les règles de la course honnête]\label{prop:g5:speed-distance-time:rules}
Deux façons honnêtes de comparer le rapide et le lent~:
\begin{enumerate}
  \item sur la \emph{même distance}, le plus rapide met \emph{moins
        de temps} --- le premier sur la ligne gagne le sprint~;
  \item en \emph{même temps}, le plus rapide parcourt \emph{plus de
        distance} --- au bout d'une minute, le nageur le plus rapide
        est tout simplement plus loin.
\end{enumerate}
Même distance ou même \omterm{def:g2:measuring-time:duration}{durée}~: fixe l'une, compare l'autre. Comparer
en changeant \emph{les deux} --- «~j'ai couru plus loin~!~» «~Oui,
mais tu as couru tout l'après-midi~!~» --- ne décide rien.
\end{proposition}

\begin{example}[Malhonnête et honnête]\label{ex:g5:speed-distance-time:fair}
Lena a fait \qty{10}{km} à vélo~; Marco en a fait \qty{15}{km}. Marco
est-il plus rapide~? Impossible de le dire --- il a peut-être roulé
trois fois plus longtemps. Mais si tous deux ont roulé exactement une
heure, la comparaison devient honnête~: même \omterm{def:g2:measuring-time:duration}{durée}, plus de distance,
Marco est plus rapide. Ce petit «~en une heure~» est sur le point de
devenir notre expression préférée.
\end{example}

\section{La vitesse~: la distance parcourue en une heure}

\begin{definition}[Vitesse]\label{def:g5:speed-distance-time:speed}
La \emph{vitesse}\index{vitesse} d'un voyageur régulier est la
distance qu'il parcourt pendant une part de temps fixée --- le plus
souvent une heure. «~Cette voiture parcourt $60$ kilomètres en une
heure~» nomme sa vitesse~; un marcheur allant bon train fait environ
$5$ kilomètres en une heure, un cycliste de course $40$. Pour les
traînards, nous prenons des parts plus commodes~: un escargot
parcourt environ $5$ \omterm{def:g2:measuring-length:units}{mètres} en une heure.
\end{definition}

\begin{omfigure}
\begin{tikzpicture}[scale=0.85]
  \draw[-Stealth, very thick] (0,0) -- (11.2,0);
  \foreach \x/\name/\v in {0.9/escargot/{$5$ m}, 3.1/marcheur/{$5$ km},
      5.3/cycliste/{$20$ km}, 7.5/voiture/{$90$ km},
      9.7/avion/{$900$ km}} {
    \fill[omDef] (\x,0) circle (2.6pt);
    \node[above, font=\footnotesize] at (\x,0.25) {\name};
    \node[below, font=\footnotesize] at (\x,-0.25) {\v};
  }
  \node[below, font=\small] at (5.6,-0.9)
    {ce que chacun parcourt en une seule heure};
\end{tikzpicture}

{\small Une échelle de \omterm{def:g5:speed-distance-time:speed}{vitesses}~: la distance que chaque voyageur
parcourt en une heure --- des cinq \omterm{def:g2:measuring-length:units}{mètres} de l'escargot aux neuf
cents kilomètres de l'avion de ligne.}
\end{omfigure}

\begin{method}[Mesurer ta propre vitesse]\label{met:g5:speed-distance-time:measure}
\begin{enumerate}
  \item \omterm{def:g2:measuring-length:measure}{Mesurer} ou trouver une longueur d'exactement \qty{100}{m}
        (beaucoup de pistes d'athlétisme en tracent une)~;
  \item la parcourir à ton allure naturelle pendant qu'un camarade
        te chronomètre --- supposons que cela prenne une minute~;
  \item une heure contient $60$ de ces minutes~: en une heure, tu
        parcourrais donc $60 \times 100 = \qty{6000}{m}$~;
  \item dit en kilomètres~: ta \omterm{def:g5:speed-distance-time:speed}{vitesse} de marche est de $6$
        kilomètres en une heure.
\end{enumerate}
Recommence en courant et compare --- la plupart des enfants doublent
à peu près.
\end{method}

\section{Prévoir les distances}

\begin{proposition}[La distance à partir de la
vitesse]\label{prop:g5:speed-distance-time:distance}
Un voyageur régulier parcourt de nouveau sa distance d'une heure à
chaque heure suivante. Donc~:
\[
\text{distance parcourue} = \text{distance en une heure} \times
\text{nombre d'heures}.
\]
Une voiture faisant $60$ kilomètres par heure parcourt, en $3$
\omterm{ex:g2:measuring-time:clock}{heures}, $60 \times 3 = \qty{180}{km}$.
\end{proposition}

\begin{example}[Calculs de trajet]\label{ex:g5:speed-distance-time:journeys}
Le marcheur à $4$ kilomètres par heure, sorti $2$ \omterm{ex:g2:measuring-time:clock}{heures}~:
$4 \times 2 = \qty{8}{km}$. Le cycliste à $15$ kilomètres par heure,
roulant $4$ \omterm{ex:g2:measuring-time:clock}{heures}~: $15 \times 4 = \qty{60}{km}$. L'avion de ligne à
$900$ kilomètres par heure, volant $5$ \omterm{ex:g2:measuring-time:clock}{heures}~:
$900 \times 5 = \qty{4500}{km}$ --- un continent traversé entre deux
repas.
\end{example}

\begin{omfigure}
\begin{tikzpicture}
\begin{axis}[omaxis, width=10.5cm, height=5.6cm,
    xlabel={heures de route}, ylabel={kilomètres parcourus},
    xtick={0,1,2,3}, ytick={0,15,30,45,60},
    ymin=0, ymax=68, xmin=0, xmax=3.4]
  \addplot[very thick, omDef] coordinates
    {(0,0) (1,15) (2,30) (3,45)};
  \addplot[very thick, omThm] coordinates
    {(0,0) (1,4) (2,8) (3,12)};
  \node[right, font=\small, omDef] at (axis cs:2.05,33)
    {cycliste~: $15$ par heure};
  \node[right, font=\small, omThm] at (axis cs:2.05,10.5)
    {marcheur~: $4$ par heure};
\end{axis}
\end{tikzpicture}

{\small Des trajets réguliers tracés en droites~: chaque heure ajoute
la même part de kilomètres. La droite du voyageur le plus rapide
monte plus fort.}
\end{omfigure}

\begin{remark}[«~Régulier~» est une simplification]\label{rem:g5:speed-distance-time:steady}
Les vrais voyageurs sont rarement réguliers~: la voiture ralentit
dans les villages, le cycliste dévale les descentes et peine dans les
montées, tu traînes devant la boulangerie. Nos calculs font comme si
tout le trajet se déroulait à une seule allure fidèle --- une
simplification utile, bien suffisante pour préparer un voyage.
Comment être honnête à propos d'un trajet à plusieurs allures --- et
ce que «~moyenne~» veut vraiment dire --- est une histoire pour une
année plus avancée, quand tu posséderas les mathématiques qu'il
faut.
\end{remark}

\begin{example}[Lire les vitesses autour de toi]\label{ex:g5:speed-distance-time:signs}
Les panneaux routiers et les tableaux de bord parlent en
kilomètres-en-une-heure --- le panneau rond indiquant $50$ dit aux
\omterm{def:g4:conductors-insulators:conductor}{conducteurs}~: parcourez au plus $50$ kilomètres en une heure, une
allure choisie pour qu'un enfant courant après un ballon puisse
encore être épargné. L'écran d'un train à grande \omterm{def:g5:speed-distance-time:speed}{vitesse} affiche
parfois fièrement $300$~; les panneaux de randonnée en montagne
préfèrent les \omterm{ex:g2:measuring-time:clock}{heures} aux kilomètres --- «~lac~: 2 h~» --- en faisant
confiance au marcheur pour connaître sa propre \omterm{def:g5:speed-distance-time:speed}{vitesse}.
\end{example}

\section{Exercices}

\begin{exercise}[$\star$]\label{exo:g5:speed-distance-time:1}
Énonce les deux façons honnêtes de comparer le rapide et le lent.
Pourquoi «~j'ai couru plus loin que toi~» ne prouve-t-il rien à lui
seul~?
\end{exercise}

\begin{exercise}[$\star$]\label{exo:g5:speed-distance-time:2}
Deux nageurs traversent le même bassin~; Ana met $30$ secondes, Bo en
met $25$. Qui est le plus rapide, et par quelle règle de la course
honnête~?
\end{exercise}

\begin{exercise}[$\star$]\label{exo:g5:speed-distance-time:3}
Au bout d'une heure, un tracteur a parcouru \qty{20}{km} et un
scooter \qty{45}{km}. Qui est le plus rapide, et par quelle règle~?
\end{exercise}

\begin{exercise}[$\star$]\label{exo:g5:speed-distance-time:4}
Que veut dire «~ce train parcourt $200$ kilomètres en une heure~»~?
Jusqu'où va-t-il en $2$ \omterm{ex:g2:measuring-time:clock}{heures}~? En une demi-heure~?
\end{exercise}

\begin{exercise}[$\star$]\label{exo:g5:speed-distance-time:5}
Range par \omterm{def:g5:speed-distance-time:speed}{vitesse}~: un marcheur ($5$ kilomètres par heure)~; \ un
pigeon ($80$ kilomètres par heure)~; \ un escargot ($5$
\emph{\omterm{def:g2:measuring-length:units}{mètres}} par heure)~; \ un bus de ville ($30$ kilomètres par
heure).
\end{exercise}

\begin{exercise}[$\star$]\label{exo:g5:speed-distance-time:6}
Tu parcours tranquillement \qty{100}{m} balisés en $2$ minutes.
Combien de \omterm{def:g2:measuring-length:units}{mètres} parcourrais-tu en une heure à cette allure~? Dis ta
\omterm{def:g5:speed-distance-time:speed}{vitesse} en kilomètres par heure.
\end{exercise}

\begin{exercise}[$\star$]\label{exo:g5:speed-distance-time:7}
Une voiture tient un régulier $80$ kilomètres par heure pendant $3$
\omterm{ex:g2:measuring-time:clock}{heures}. Quelle distance parcourt-elle~? Et un cycliste à $15$
kilomètres par heure pendant les mêmes $3$ \omterm{ex:g2:measuring-time:clock}{heures}~?
\end{exercise}

\begin{exercise}[$\star$]\label{exo:g5:speed-distance-time:8}
Sur le graphique des trajets du chapitre, comment repérer d'un coup
d'œil le voyageur le plus rapide~? Lis dessus~: quelle distance le
cycliste a-t-il parcourue après $2$ \omterm{ex:g2:measuring-time:clock}{heures}~?
\end{exercise}

\begin{exercise}[$\star\star$]\label{exo:g5:speed-distance-time:9}
Un ferry parcourt $30$ kilomètres en une heure. La traversée fait
\qty{90}{km}. En utilisant
\cref{prop:g5:speed-distance-time:distance} à l'envers --- essaie
heure après heure --- combien de temps dure la traversée~?
\end{exercise}

\begin{exercise}[$\star\star$]\label{exo:g5:speed-distance-time:10}
Deux villes sont distantes de \qty{60}{km}. Mara part de l'une à vélo
à $20$ kilomètres par heure~; au même instant, Jon part de l'autre
vers elle à $10$ kilomètres par heure. De combien de kilomètres
l'écart entre eux diminue-t-il chaque heure~? Au bout de combien
d'\omterm{ex:g2:measuring-time:clock}{heures} se rencontrent-ils~?
\end{exercise}

\begin{exercise}[$\star\star$]\label{exo:g5:speed-distance-time:11}
La vieille fable, avec des nombres~: la course fait \qty{120}{m} de
long. Le lièvre court $60$ \omterm{def:g2:measuring-length:units}{mètres} par minute --- mais après une
minute de course, il s'arrête pour une sieste de $70$ minutes, puis
repart. La tortue avance de $2$ \omterm{def:g2:measuring-length:units}{mètres} par minute et ne s'arrête
jamais. Où est le lièvre quand il commence sa sieste~? À quel instant
chacun franchit-il la ligne d'arrivée~? Qui gagne~?
\end{exercise}
